\section{The HNF Algorithm and Formalization}
\label{sec:algorithm}

\begin{figure}[t]
  \centering
  \begin{tikzpicture}[x=1cm,y=1cm,font=\small]
  \node[fig/node plain, minimum width=1.45cm, minimum height=1.35cm] (input) at (0.85,1.55)
    {\textit{Input}\\$A$};
  \node[fig/node, minimum width=2.40cm, minimum height=1.35cm, right=0.40cm of input] (clear)
    {{\ttfamily clearFirstColumn}\\[-0.15ex]{\footnotesize pivot; zeros below}};
  \node[fig/node, minimum width=2.35cm, minimum height=1.35cm, right=0.40cm of clear] (recurse)
    {{\ttfamily recurse}\\[-0.15ex]{\footnotesize on $\lowerRight$}};
  \node[fig/node, minimum width=1.85cm, minimum height=1.35cm, right=0.40cm of recurse] (lift)
    {{\ttfamily diagLift}\\[-0.15ex]{\footnotesize embed $U$}};
  \node[fig/node, minimum width=2.40cm, minimum height=1.35cm, right=0.40cm of lift] (reduce)
    {{\ttfamily reduceTopRow}\\[-0.15ex]{\footnotesize canonical remainders}};
  \node[fig/node plain, minimum width=1.45cm, minimum height=1.35cm, right=0.40cm of reduce] (output)
    {\textit{Output}\\$H$};

  \draw[fig/arrow] (input) -- (clear);
  \draw[fig/arrow] (clear) -- (recurse);
  \draw[fig/arrow] (recurse) -- (lift);
  \draw[fig/arrow] (lift) -- (reduce);
  \draw[fig/arrow] (reduce) -- (output);

  \node[fig/stage] at ($(clear.north)+(0,0.55)$) {Stage 1};
  \node[fig/stage] at ($(recurse.north)!0.5!(lift.north)+(0,0.55)$) {Stage 2};
  \node[fig/stage] at ($(reduce.north)+(0,0.55)$) {Stage 3};

  \node[fig/banner, text width=12.6cm, font=\footnotesize] at ($(lift.south)+(0,-1.05)$)
    {global invariant: unimodular $U$ and $U \cdot A = B$};
\end{tikzpicture}

  \caption{The three-stage structure of the recursive HNF kernel. A single
  unimodular left certificate carries the invariant \(U \cdot A = B\)
  throughout.}
  \label{fig:hnf-pipeline}
\end{figure}

The executable kernel is \code{hermiteNormalFormFin}. It returns an internal
result package containing the left multiplier \(U\), the output matrix \(H\), the
equation \(U A = H\), a chosen inverse \(U^{-1}\), both multiplication
identities, and a proof that \(H\) is in
\code{IsHermiteNormalFormFin}. Determinant-based unimodularity is a derived
view. The recursion follows the three-stage pattern shown in
Figure~\ref{fig:hnf-pipeline}.

\subsection{Stage 1: clear the first column}

For a nonempty matrix, the kernel first checks whether the entire first column is
zero. If so, it recurses on \(\tailCols A\). Otherwise it finds the first nonzero
entry of the first column, swaps that row into the top position, and normalizes
the new pivot with \code{normUnit}. It then clears the rest of the first column
by iterating \code{clearFirstColumnStep}.

\begin{lstlisting}[style=leanstyle,caption={The local first-column elimination step.},label={lst:clearstep}]
def clearFirstColumnStep (i : Fin (m + 1)) (t : LeftTransform A) :
    LeftTransform A :=
  if ComputableEuclideanOps.isZeroB (t.B i.succ 0) = true then
    t
  else
    let tSwap := t.trans (LeftTransform.swap t.B i.succ (1 : Fin (m + 2)))
    let tBez := tSwap.trans (LeftTransform.topBezout tSwap.B)
    tBez.trans
      (LeftTransform.unitSmul tBez.B 0
        (ComputableEuclideanOps.normUnitUnit (tBez.B 0 0)))
\end{lstlisting}

This is the canonical micro-step of the algorithm. If the current entry in the
first column is already zero, nothing happens. Otherwise the target row is moved
into position \(1\), the first two rows are rewritten by a lifted Bezout matrix,
and the new pivot is normalized once again. Repeating this step down the column
accumulates the gcd of all processed first-column entries in the top-left corner
and forces zeros below it.

The underlying \(2 \times 2\) gadget lives in
\module{NormalForms.Matrix.Elementary.Basic}. For \(a,b \in R\), the matrix
\code{bezoutReductionMatrix a b}, abbreviated here by \(B(a,b)\), satisfies
\[
  B(a,b) \cdot (a,b)^\mathsf{T}
  = (g, 0)^\mathsf{T},
\]
where \(g\) is the normalized result of the constructive extended gcd. The
same primitive constructs its explicit \(2\times2\) inverse and proves both
matrix identities; no determinant or matrix inversion is used to recover it.
Dedicated matrix-vector and semantic lemmas live in the elementary layer. The
HNF layer lifts this gadget to larger
matrices by \code{topBezoutMatrix}, whose tail block is the identity. The lemmas
\code{topBezoutMatrix_mul_topLeft} and
\code{topBezoutMatrix_mul_rowOneFirstCol} transport the gcd computation and the
new zero to the leading column of the full matrix, while
\code{Bezout_preserves_zeros} shows that rows below the active \(2 \times 2\)
window keep their first-column zeros.

\subsection{Stage 2: recurse on the lower-right block}

Once the first column has the desired pivot shape, the kernel recurses on
\(\lowerRight A\). The result is a left certificate for the smaller matrix. To
bring it back to the original dimension, the development uses
\code{diagLiftMatrix U}, the block matrix
\[
  \begin{bmatrix}
    1 & 0 \\
    0 & U
  \end{bmatrix}.
\]

The corresponding constructor \code{LeftTransform.diagLift} is the key example
of why the recursive state is a \code{LeftTransform} and not merely a row log.
Its fundamental preservation lemmas are:

\begin{itemize}
  \item \code{diagLiftMatrix_mul_topRow}: the top row is unchanged;
  \item \code{lowerRight_diagLiftMatrix_mul}: the lower-right block is updated
    exactly by multiplication with the recursive certificate;
  \item the block-lift inverse certificate: the lifted forward and backward
    matrices satisfy both inverse identities.
\end{itemize}

These three facts are enough to splice the recursive call back into the ambient
matrix without reopening any log semantics. The algorithm composes the current
certificate with the lifted recursive certificate by \code{LeftTransform.trans}.

\subsection{Stage 3: reduce the top row}

After the lower-right block is itself in HNF, the remaining work is to reduce the
top row modulo the pivots that now appear in lower rows. This is the purpose of
\code{reduceTopRowLoop}. For a lower row \(i+1\), the algorithm searches for the
first nonzero entry in the tail of that row. If it occurs at \(j+1\), it adds a
multiple of row \(i+1\) to the top row with coefficient
\[
  -\,\frac{A_{0,j+1}}{A_{i+1,j+1}}.
\]
In code this is \code{topReductionCoeff}. The lemma
\code{applyRowOperation_reduceTop_pivot} proves that the resulting top entry is
literally \(A_{0,j+1} \% A_{i+1,j+1}\). The companion lemma
\code{reduceTopRowStep_current_reduced} then closes the reducedness proof by a
single call to \code{CanonicalMod.mod_mod}.

Just as importantly, the reduction loop has carefully controlled side effects.
The lemmas \code{reduceTopRowLoop_lowerRight},
\code{reduceTopRowLoop_lowerRow}, \code{reduceTopRowLoop_below_zero}, and
\code{reduceTopRowLoop_topLeft} show that lower rows are unchanged, the
lower-right block is unchanged, zeros below the top-left pivot are preserved, and
the pivot itself is not disturbed. The final theorem
\code{reduceTopRowLoop_reduced_prefix} packages these preservation facts into
the exact reduction condition needed by the \code{pivot} constructor of
\code{IsHermiteNormalFormFin}.

\subsection{Public extraction and current execution boundary}

The public functions return strong results directly. Callers may use
\code{hermiteNormalFormFin}, provide a \code{MatrixIndexing}, supply explicit
row and column equivalences, request increasing order, or choose the
noncomputable convenience wrapper \code{hermiteNormalForm}. The stored inverse
certificate yields both matrix identities and the compatibility theorem
\code{hermiteNormalForm_unimodular}; \code{hermiteNormalForm_isHermite} proves
normality relative to the stored indexing.

The repository includes exact-output \code{\#eval} harnesses over \code{Int}
and \code{Polynomial Rat}, plus a dedicated native polynomial certificate
executable. The regression corpus covers signs, units, zero dimensions,
rectangles, rank deficiency, pivot improvement, transformation equations, and
every stored inverse identity. Rational-polynomial arithmetic, long division,
normalization, and extended gcd are direct executable definitions connected
to \mathlib{} semantics by proofs. Smith stabilization uses
\code{ComputableEuclideanOps.measure} and three strict-descent theorems rather
than enumerating normalized factors.
